\section{E1: effective noise and fixed-factor training}
\label{sec:e1}
\subsection{Why the effective matrix matters}
Condition on the signal-updated factors and the just-updated user matrix.
Let $Z_A$ and $Z_B$ have independent, identically distributed Gaussian entries
with SD $\tau=\eta_s\sigma C/(qN)$. The immediate effective-matrix shock is
\begin{equation}
 D=(A+Z_A)(B+Z_B)-AB=AZ_B+Z_AB+Z_AZ_B.
 \label{eq:productshock}
\end{equation}
The three terms have zero mean and zero expected pairwise Frobenius inner
products, but they do not all need to be independent. Their expected energies
give
\begin{align}
 \E\norm{D}_F^2
 &=\tau^2\bigl(M\norm{B}_F^2+d\norm{A}_F^2\bigr)+Mdr\tau^4,
 \label{eq:qmoment}\\
 \E\norm{PD^T}_F^2
 &=\tau^2\bigl(M\norm{PB^T}_F^2+\norm{A}_F^2\norm{P}_F^2\bigr)
   +Mr\tau^4\norm{P}_F^2.
 \label{eq:smoment}
\end{align}
Thus coordinate dimension is only one determinant of prediction disturbance.
A method with modest effective item noise can still have a large score shock
if its local user vectors grow. Conversely, multiplying all of a user's
scores by a positive constant changes their scale without changing their
ranking. Norms, score RMS, margins, and ranking outcomes must be interpreted
together. These conditional calculations do not describe the full distribution
after many rounds of state-dependent learning.

FixedB removes the two terms involving a noisy $B$. With $BB^T=I_r$,
$D=Z_AB$ and $\norm{D}_F=\norm{Z_A}_F$. This is a clean linear diagnostic
baseline. It also changes clipping geometry and aggregation behavior, so
better conditional noise moments do not guarantee better trained utility.
Eight synthetic Monte Carlo checks corroborate the implemented moment
calculations; derivations and assumptions are in \cref{app:theory}.

\subsection{Controlled utility comparison}
% Generated by scripts/generate_tables.py; do not edit numeric cells.
\begin{table}[tbp]\centering\small
\caption{E1 test utility with ranks selected on validation at each privacy level.}\label{tab:e1}
\setlength{\tabcolsep}{4pt}
\begin{tabular}{lrrrr}\toprule
$\epsilon$ & Full & Two & FixedB & DP popularity \\
\midrule
8 & 0.052271 & 0.051158 (r32) & 0.048978 (r32) & 0.049548 \\
4 & 0.042169 & 0.045807 (r32) & 0.038959 (r32) & 0.049428 \\
2 & 0.025428 & 0.040841 (r8) & 0.025286 (r16) & 0.049507 \\
1 & 0.011461 & 0.025136 (r4) & 0.011958 (r16) & 0.048389 \\
\bottomrule\end{tabular}
\tabnote{Five-seed NDCG@10 means. The primary E1 family compares matched ranks, not these separately selected ranks. Source: test\_means.csv and rank\_selection.csv.}
\end{table}

The primary family compares FixedB with Two at each matched rank and privacy
level. None of the sixteen adjusted comparisons favors FixedB. Eight favor
Two and eight include zero. At target two, FixedB is worse at all four ranks.
The rank-8 difference is $-0.018662$, with adjusted interval
$[-0.027695,-0.010368]$. This is stronger evidence against the proposed
standalone fixed-factor improvement than an informal comparison of two
selected example runs.

The validation-selected presentation in \cref{tab:e1} leads to the same
practical conclusion, while answering a different question. Two chooses
ranks 32/32/8/4 at privacy targets 8/4/2/1; FixedB chooses 32/32/16/16.
Two reaches $0.040841$ at target two and $0.025136$ at target one, whereas
FixedB reaches $0.025286$ and $0.011958$. These selected-rank curves should
not be described as the matched-rank primary statistical family.

FixedB can perform well in no-noise controls, and it often has smaller
immediate effective noise. Its poorer private utility therefore does not
show that the implementation lacks all ranking capacity. It shows that
removing one noise source, with these transferred clipping and search choices,
does not overcome the remaining optimization and learning limitations.
The grid boundary, parameterization-dependent clipping, and different selected
server rates remain qualifications on any broader statement about optimal
fixed-factor performance.

\reportfigure{03_effective_noise_utility.pdf}{Controlled E1 utility}
{Full, Two, FixedB, and private popularity test NDCG@10 with five-seed SD.
Ranks for Two and FixedB are selected separately on validation at each level.
The primary tests instead compare matched ranks. Collaboration uses 50
rounds and popularity one separately calibrated bounded release. Every
collaborative point mean at target epsilon at most four is below private
popularity; this is not a family-wide significance claim against popularity.}
{fig:e1}

\subsection{The private popularity baseline changes the reference point}
The bounded no-noise popularity model obtains $0.049885$ test NDCG@10.
Its private mean at target one is $0.048389\pm0.001125$, retaining about
97\% of that bounded reference. The corresponding means at targets 8, 4,
and 2 are $0.049548$, $0.049428$, and $0.049507$. All E1 collaborative
point means at targets at most four lie below the corresponding private
popularity mean. This is a descriptive statement covering the retained
settings, not a multiplicity-adjusted superiority theorem.

The comparison is useful even though the mechanisms allocate their budgets
differently. Both protect an entire user's bounded contribution at the
reported target, but popularity uses one full-population release and does
not attempt personalized collaborative representation learning. Its strength
suggests that the useful private signal in these settings is not exhausted
by the iterative methods' current optimization. It does not prove that
personalization is fundamentally impossible under the same privacy target.

\subsection{Measured geometry, clipping, and numerical behavior}
\reportfigure{04_noise_and_clipping.pdf}{Effective noise, score scale, and clipping}
{E1 diagnostics at target epsilon one. Left: mean over seeds of each run's
root-mean-squared Frobenius norm of the immediate effective item shock.
Center: mean score-shock RMSE, on a logarithmic axis retaining the large finite
Two-r32 outlier. Right: selected-client clipping fraction. These aggregates
are diagnostics, not direct causal measures of NDCG loss. Source:
\texttt{geometry\_means.csv}; per-run evidence remains available.}{fig:geometry}
At target one, the Full effective-noise norm is approximately $10.369$.
FixedB ranks 4/8/16/32 have norms $2.592$, $3.667$, $5.186$, and $7.333$;
the corresponding Two values are approximately $3.602$, $4.817$, $6.542$,
and $33.788$. These are averages of per-run RMS Frobenius norms over rounds,
not per-entry noise standard deviations. Different selected server rates
also affect the comparison.

The expected bilinear term accounts for only about $0.2$--$0.3$\% of total
effective-matrix noise energy in the selected Two settings at target one.
Its presence is algebraically important, but its measured energy does not
support calling it the dominant cause of the observed utility failures.
First-order terms, factor scales, clipping, and local learning remain relevant.

Clipping at the same coordinate threshold affects the methods differently.
At target one, approximately $54.1$\% of selected Full updates are clipped;
FixedB fractions range from $25.6$\% to $46.9$\%; Two fractions range from
$54.6$\% to $82.5$\%. These fractions accompany different effective signal
steps, not just different noise counts. Equal $C$ across parameterizations
does not mean equal effective-$Q$ sensitivity or equal retained signal.

Some Two trajectories have very large but finite local user norms. The mean
score-shock RMSE for Two-r32 at target one is approximately
$9.846\times10^{16}$, dominated by extreme state scale. An original float32
norm summary overflowed even though stored vectors were finite; analysis
was corrected by reading those vectors in float64. Original logs, training,
and held-out ranks were preserved. The large values are reported rather
than excluded, but a mean dominated by such an outlier is not a typical-user
effect. Ranking diagnostics and seed-level distributions are necessary
companions to the RMS summary.

\subsection{Activity groups and personalization ablations}
At target two, validation-selected Two-r8 obtains test NDCG@10 of
$0.048722$, $0.027027$, and $0.046489$ in the low-, medium-, and
high-activity groups. Private popularity gives $0.055746$, $0.036872$,
and $0.055662$. Thus the aggregate gap in this comparison is not explained
solely by one activity group. These are descriptive group means, without
an adjusted subgroup test. The cold-target group is too small for a reliable
comparison: Two-r8 and popularity both score zero there at target two,
while Full has a small nonzero value from very few top-ten outcomes.

Local-vector ablations expose a further distinction between storing a user
vector and obtaining useful personalization from it. In the Full no-DP
control, validation NDCG is $0.053136$ using the user's own vector,
$0.042725$ with permuted vectors, and $0.047421$ with the mean vector.
At target one the corresponding values are $0.008883$, $0.011147$,
and $0.030485$. FixedB-r8 likewise has little own-versus-permuted difference
at target one ($0.008580$ versus $0.008481$), while its mean-vector
diagnostic is $0.018912$. These contrasts show that local parameter
persistence alone does not guarantee a personalization benefit under noise.

The mean-vector calculation uses the collection of private local states
and is evaluated retrospectively. It is not a free public statistic,
a deployable private method, or a new selected baseline. There is no
confirmatory ablation significance claim. Its role is to motivate further
investigation of how local user learning interacts with the noisy item
representation, which E2 studies through explicit interventions.

\subsection{Recommendation-specific implications}
The item factor is strongly rectangular: $A$ has $1682r$ entries and $B$
only $64r$. Fixing $B$ therefore saves $64/(1682+64)=3.67\%$ of repeatedly
shared coordinates relative to Two at the same rank, not half the payload.
At rank 8 the FixedB payload is 107,648 bytes versus Two's 111,744. Most
savings relative to Full come from low rank, which both methods already use.

The activity, cold-target, and local-vector ablation results are retained in
the supplement. They provide descriptive information about heterogeneous
personalization, but they do not establish a subgroup-specific improvement.
In particular, a random public basis supplies no semantic information for
an unobserved item. A metadata-only personalizer and a small private residual
remain distinct future comparisons. E1's substantive outcome is the failure
of a plausible fixed-factor improvement under controlled testing, together
with a clearer set of measurements for investigating why.
\FloatBarrier
